\documentclass[12pt,a4paper]{article}
\usepackage[utf8]{inputenc}
\usepackage[T1]{fontenc}
\usepackage[french]{babel}
\usepackage{geometry}
\geometry{margin=2.2cm}

\title{Documentation des Sprints - VMS}
\author{}
\date{}

\begin{document}
\maketitle

\section*{Sprint 1 - Authentification et gestion des rôles}
\subsection*{Explication}
Cette fonctionnalité met en place l’authentification sécurisée des utilisateurs (employé/visiteur) et la gestion des autorisations par rôle. L’objectif métier est de garantir que chaque action sensible (consultation, modification, validation d’accès) soit exécutée uniquement par les profils autorisés, tout en offrant une connexion fluide pour les utilisateurs mobiles et les administrateurs.

\subsection*{Diagramme de classes}
\begin{verbatim}
classDiagram
%% Sprint 1
class AuthController {
  -jwtService: JwtService
  -employeeService: EmployeeService
  +loginEmployee(email: String, password: String): String
  +loginVisitor(email: String, totpCode: String): String
}

class EmployeeService {
  -employeeRepository: EmployeeRepository
  -passwordEncoder: PasswordEncoder
  +authenticate(email: String, password: String): Employee
  +loadUserByUsername(email: String): Employee
}

class JwtService {
  -privateKey: String
  -publicKey: String
  -ttlSeconds: long
  +generateToken(subject: String, scope: String): String
  +validateToken(token: String): boolean
}

class Employee {
  -email: String
  -passwordHash: String
  -role: String
  +checkPassword(raw: String): boolean
  +toScope(): String
}

class AuthServiceMobile {
  -dioClient: DioClient
  -tokenHolder: TokenHolder
  +loginEmployee(email: String, password: String): String
  +saveToken(jwt: String): void
}

class TokenHolder {
  -accessToken: String
  -refreshAt: String
  -userEmail: String
  +getToken(): String
  +clearToken(): void
}

AuthController "1" --> "1" EmployeeService : dependance
AuthController "1" --> "1" JwtService : dependance
EmployeeService "1" --> "0..*" Employee : utilise
EmployeeService "1" --> "1" EmployeeRepository : dependance
AuthServiceMobile "1" --> "1" TokenHolder : association
AuthServiceMobile "1" --> "1" AuthController : appel API
\end{verbatim}

\subsection*{Diagramme de cas d’utilisation}
\begin{verbatim}
usecaseDiagram
%% Sprint 1
actor "Utilisateur" as U
actor "Admin" as A

rectangle "VMS - Authentification" {
  usecase "Se connecter" as UC1
  usecase "Recevoir JWT" as UC2
  usecase "Acceder aux ecrans autorises" as UC3
  usecase "Gerer les roles et permissions" as UC4
  usecase "Reinitialiser les acces" as UC5
}

U --> UC1
U --> UC2
U --> UC3
A --> UC4
A --> UC5
A --> UC3
\end{verbatim}

\subsection*{Diagramme de séquence}
\begin{verbatim}
sequenceDiagram
%% Sprint 1
participant Mobile as App Flutter
participant API as API /vms
participant C as AuthController
participant S as EmployeeService
participant R as EmployeeRepository
participant DB as PostgreSQL

Mobile->>API: POST /employees/token (Basic)
API->>C: loginEmployee(email,password)
C->>S: authenticate(email,password)
S->>R: findByEmail(email)
R->>DB: SELECT employee by email
DB-->>R: employee row
R-->>S: Employee
S-->>C: Employee valide
C-->>API: JWT
API-->>Mobile: 200 token
Mobile->>Mobile: saveToken(token)
\end{verbatim}

\section*{Sprint 2 - Gestion des réunions}
\subsection*{Explication}
Cette fonctionnalité couvre le cycle de vie des réunions : création, consultation, modification et transitions de statut. L’objectif métier est de centraliser l’organisation des visites et de synchroniser les informations entre organisateur, participants et système de contrôle d’accès.

\subsection*{Diagramme de classes}
\begin{verbatim}
classDiagram
%% Sprint 2
class MeetingController {
  -meetingService: MeetingService
  -securityContext: SecurityContext
  +createMeeting(dto: MeetingDto): MeetingDto
  +updateMeeting(id: Long, dto: MeetingDto): MeetingDto
}

class MeetingService {
  -meetingRepository: MeetingRepository
  -officeRepository: OfficeRepository
  -employeeRepository: EmployeeRepository
  +create(dto: MeetingDto, actor: String): Meeting
  +changeStatus(id: Long, status: String): void
}

class MeetingRepository {
  -entityManager: EntityManager
  -tableName: String
  +save(meeting: Meeting): Meeting
  +findById(id: Long): Meeting
}

class Meeting {
  -id: Long
  -topic: String
  -status: String
  -startTime: String
  -endTime: String
  +isEditable(): boolean
  +setStatus(newStatus: String): void
}

class MeetingServiceMobile {
  -dioClient: DioClient
  -baseUrl: String
  +fetchMeetings(page: int): List~MeetingModel~
  +createMeeting(data: Map): MeetingModel
}

class MeetingModel {
  -id: int
  -topic: String
  -status: String
  -startTime: String
  +fromJson(json: Map): MeetingModel
  +toJson(): Map
}

MeetingController "1" --> "1" MeetingService : dependance
MeetingService "1" --> "1" MeetingRepository : dependance
MeetingRepository "1" --> "0..*" Meeting : persistance
MeetingServiceMobile "1" --> "1" MeetingController : appel API
MeetingServiceMobile "1" --> "0..*" MeetingModel : conversion
\end{verbatim}

\subsection*{Diagramme de cas d’utilisation}
\begin{verbatim}
usecaseDiagram
%% Sprint 2
actor "Utilisateur" as U
actor "Admin" as A

rectangle "VMS - Reunions" {
  usecase "Creer une reunion" as UC1
  usecase "Consulter les reunions" as UC2
  usecase "Modifier une reunion" as UC3
  usecase "Changer le statut de reunion" as UC4
  usecase "Supprimer une reunion" as UC5
}

U --> UC1
U --> UC2
U --> UC3
U --> UC4
A --> UC2
A --> UC4
A --> UC5
\end{verbatim}

\subsection*{Diagramme de séquence}
\begin{verbatim}
sequenceDiagram
%% Sprint 2
participant Mobile as App Flutter
participant API as API /vms
participant C as MeetingController
participant S as MeetingService
participant R as MeetingRepository
participant DB as PostgreSQL

Mobile->>API: POST /meetings (Bearer + JSON)
API->>C: createMeeting(dto)
C->>S: create(dto, actor)
S->>R: save(meeting)
R->>DB: INSERT meeting
DB-->>R: meeting_id
R-->>S: Meeting persisted
S-->>C: MeetingDto
C-->>API: 201 Created
API-->>Mobile: MeetingDto JSON
\end{verbatim}

\section*{Sprint 3 - Génération et validation des QR codes}
\subsection*{Explication}
Cette fonctionnalité permet de générer des QR codes d’accès liés à une réunion et de les valider au scan. L’objectif métier est de sécuriser le contrôle d’entrée, réduire le temps de vérification et garantir la traçabilité des accès autorisés/refusés.

\subsection*{Diagramme de classes}
\begin{verbatim}
classDiagram
%% Sprint 3
class ParticipantController {
  -meetingService: MeetingService
  -participantService: ParticipantService
  +getSecret(meetingId: Long): UserSecretDto
  +resolveMeeting(meetingId: Long, participant: String): MeetingResolveDto
}

class ParticipantService {
  -meetingRepository: MeetingRepository
  -employeeRepository: EmployeeRepository
  -visitorRepository: VisitorRepository
  +fetchParticipantSecret(meetingId: Long, user: String): String
  +isParticipantAllowed(meetingId: Long, email: String): boolean
}

class UserSecretDto {
  -displayName: String
  -secret: String
  -meetingId: Long
  +maskSecret(): String
  +toJson(): Map
}

class QrPopupControllerMobile {
  -participantServiceMobile: ParticipantServiceMobile
  -jwtHelper: JwtHelper
  +buildQrPayload(meetingId: int, participant: String): String
  +generateQrToken(payload: Map): String
}

class MeetingScannerControllerMobile {
  -meetingServiceMobile: MeetingServiceMobile
  -jwtHelper: JwtHelper
  +scanAndDecode(raw: String): Map
  +verifyAccess(meetingId: int, participant: String): bool
}

class ParticipantServiceMobile {
  -dioClient: DioClient
  -baseUrl: String
  -tokenHolder: TokenHolder
  +getParticipantSecret(meetingId: int): Map
  +resolveMeeting(meetingId: int, participant: String): Map
}

ParticipantController "1" --> "1" ParticipantService : dependance
ParticipantService "1" --> "0..*" UserSecretDto : construit
QrPopupControllerMobile "1" --> "1" ParticipantServiceMobile : dependance
MeetingScannerControllerMobile "1" --> "1" ParticipantServiceMobile : dependance
ParticipantServiceMobile "1" --> "1" ParticipantController : appel API
\end{verbatim}

\subsection*{Diagramme de cas d’utilisation}
\begin{verbatim}
usecaseDiagram
%% Sprint 3
actor "Utilisateur" as U
actor "Admin" as A

rectangle "VMS - QR Access" {
  usecase "Generer QR d'acces" as UC1
  usecase "Scanner QR" as UC2
  usecase "Verifier validite d'acces" as UC3
  usecase "Afficher resultat autorise/refuse" as UC4
  usecase "Auditer les acces" as UC5
}

U --> UC1
U --> UC2
U --> UC3
U --> UC4
A --> UC5
A --> UC3
\end{verbatim}

\subsection*{Diagramme de séquence}
\begin{verbatim}
sequenceDiagram
%% Sprint 3
participant Scanner as App Scanner Flutter
participant API as API /vms
participant C as ParticipantController
participant S as ParticipantService
participant R as MeetingRepository
participant DB as PostgreSQL

Scanner->>API: GET /meetings/{id}/resolve?participant=email
API->>C: resolveMeeting(id, participant)
C->>S: isParticipantAllowed(id, email)
S->>R: findMeetingWithParticipants(id)
R->>DB: SELECT meeting + participants
DB-->>R: meeting data
R-->>S: Meeting aggregate
S-->>C: allowed + secret
C-->>API: MeetingResolveDto
API-->>Scanner: 200 access decision
Scanner->>Scanner: show Pass/Reject screen
\end{verbatim}

\section*{Sprint 4 - Gestion visiteurs et check-in}
\subsection*{Explication}
Cette fonctionnalité gère l’enregistrement des visiteurs, le redeem des invitations et le check-in à l’arrivée. L’objectif métier est de fluidifier l’accueil, fiabiliser les informations des visiteurs et connecter l’identité du visiteur à une réunion valide.

\subsection*{Diagramme de classes}
\begin{verbatim}
classDiagram
%% Sprint 4
class VisitorController {
  -visitorService: VisitorService
  -tokenService: TokenService
  +createVisitor(dto: VisitorDto): VisitorDto
  +checkIn(dto: CheckInDto): void
}

class VisitorService {
  -visitorRepository: VisitorRepository
  -meetingRepository: MeetingRepository
  -tokenRepository: TokenRepository
  +redeem(code: String, visitor: String): FullVisitorDto
  +checkIn(dto: CheckInDto): boolean
}

class Visitor {
  -email: String
  -firstname: String
  -lastname: String
  -company: String
  +fullName(): String
  +updateProfile(company: String, phone: String): void
}

class VisitorCheckinControllerMobile {
  -visitorServiceMobile: VisitorServiceMobile
  -userServiceMobile: UserServiceMobile
  +submitCheckIn(form: Map): bool
  +routeByUserType(email: String): String
}

class VisitorServiceMobile {
  -dioClient: DioClient
  -baseUrl: String
  -tokenHolder: TokenHolder
  +checkIn(data: Map): bool
  +redeemInvite(code: String, visitor: String): Map
}

class FullVisitorModel {
  -email: String
  -issuer: String
  -secret: String
  +fromJson(json: Map): FullVisitorModel
  +toQrLoginPayload(): Map
}

VisitorController "1" --> "1" VisitorService : dependance
VisitorService "1" --> "0..*" Visitor : manipule
VisitorCheckinControllerMobile "1" --> "1" VisitorServiceMobile : dependance
VisitorServiceMobile "1" --> "1" VisitorController : appel API
VisitorServiceMobile "1" --> "0..*" FullVisitorModel : conversion
\end{verbatim}

\subsection*{Diagramme de cas d’utilisation}
\begin{verbatim}
usecaseDiagram
%% Sprint 4
actor "Utilisateur" as U
actor "Admin" as A

rectangle "VMS - Visiteurs" {
  usecase "Enregistrer un visiteur" as UC1
  usecase "Redeem code d'invitation" as UC2
  usecase "Effectuer check-in" as UC3
  usecase "Mettre a jour profil visiteur" as UC4
  usecase "Superviser le registre visiteurs" as UC5
}

U --> UC1
U --> UC2
U --> UC3
U --> UC4
A --> UC5
A --> UC4
\end{verbatim}

\subsection*{Diagramme de séquence}
\begin{verbatim}
sequenceDiagram
%% Sprint 4
participant Mobile as App Visitor Flutter
participant API as API /vms
participant C as VisitorController
participant S as VisitorService
participant R as VisitorRepository
participant DB as PostgreSQL

Mobile->>API: POST /visitors/checkin (CheckInDto)
API->>C: checkIn(dto)
C->>S: checkIn(dto)
S->>R: findByEmail(dto.email)
R->>DB: SELECT visitor by email
DB-->>R: visitor row
R-->>S: Visitor
S->>DB: UPDATE checkin state
DB-->>S: ok
S-->>C: true
C-->>API: 200 OK
API-->>Mobile: check-in success
\end{verbatim}

\section*{Sprint 5 - Administration, offices et gouvernance opérationnelle}
\subsection*{Explication}
Cette fonctionnalité couvre les opérations d’administration (gestion offices, utilisateurs, supervision des données) pour garantir un fonctionnement stable en production. L’objectif métier est de donner à l’admin une vision de contrôle globale, avec des actions de maintenance sécurisées et traçables.

\subsection*{Diagramme de classes}
\begin{verbatim}
classDiagram
%% Sprint 5
class OfficeController {
  -officeService: OfficeService
  -employeeService: EmployeeService
  +createOffice(dto: OfficeDto): OfficeDto
  +updateOffice(id: Long, dto: OfficeDto): OfficeDto
}

class OfficeService {
  -officeRepository: OfficeRepository
  -employeeRepository: EmployeeRepository
  -auditService: AuditService
  +assignManager(officeId: Long, managerEmail: String): void
  +deleteOffice(officeId: Long): void
}

class Office {
  -id: Long
  -name: String
  -city: String
  -country: String
  +changeLocation(city: String, country: String): void
  +setManager(email: String): void
}

class AdminDashboardControllerMobile {
  -officeServiceMobile: OfficeServiceMobile
  -employeeServiceMobile: EmployeeServiceMobile
  +loadAdminOverview(): Map
  +triggerOfficeUpdate(data: Map): bool
}

class OfficeServiceMobile {
  -dioClient: DioClient
  -baseUrl: String
  -tokenHolder: TokenHolder
  +fetchOffices(): List~OfficeModel~
  +updateOffice(id: int, payload: Map): OfficeModel
}

class OfficeModel {
  -id: int
  -name: String
  -city: String
  -country: String
  +fromJson(json: Map): OfficeModel
  +toJson(): Map
}

OfficeController "1" --> "1" OfficeService : dependance
OfficeService "1" --> "0..*" Office : manipule
AdminDashboardControllerMobile "1" --> "1" OfficeServiceMobile : dependance
OfficeServiceMobile "1" --> "1" OfficeController : appel API
OfficeServiceMobile "1" --> "0..*" OfficeModel : conversion
\end{verbatim}

\subsection*{Diagramme de cas d’utilisation}
\begin{verbatim}
usecaseDiagram
%% Sprint 5
actor "Utilisateur" as U
actor "Admin" as A

rectangle "VMS - Administration" {
  usecase "Consulter offices" as UC1
  usecase "Creer / modifier office" as UC2
  usecase "Assigner manager office" as UC3
  usecase "Supprimer entites obsoletes" as UC4
  usecase "Superviser coherence des donnees" as UC5
}

U --> UC1
A --> UC1
A --> UC2
A --> UC3
A --> UC4
A --> UC5
\end{verbatim}

\subsection*{Diagramme de séquence}
\begin{verbatim}
sequenceDiagram
%% Sprint 5
participant Mobile as Admin Mobile App
participant API as API /vms
participant C as OfficeController
participant S as OfficeService
participant R as OfficeRepository
participant DB as PostgreSQL

Mobile->>API: PATCH /offices/{id} (OfficeDto)
API->>C: updateOffice(id, dto)
C->>S: assignManager(id, dto.manager)
S->>R: findById(id)
R->>DB: SELECT office by id
DB-->>R: office row
R-->>S: Office
S->>R: save(office updated)
R->>DB: UPDATE office
DB-->>R: ok
R-->>S: Office updated
S-->>C: OfficeDto
C-->>API: 200 updated
API-->>Mobile: OfficeDto JSON
\end{verbatim}

\end{document}
